\chapter{Discussion}

\section{What the Current Results Suggest}

RQ1 begins with a descriptive prevalence premise and then uses two complementary interventions to study operational-information routing. In a frozen pool of 29,292 public primary skill documents, each complete \texttt{SKILL.original.md} file---including its \texttt{name}, \texttt{description}, and Markdown content---received a model-assisted binary presence review. All eight candidate information types recur frequently: from 64.91\% for few-shot examples to 83.34\% for workflow/procedure; use condition, output/artifact, dependency/resource, and success/verification occur in 80.68\%, 79.76\%, 75.32\%, and 71.58\% of documents respectively. This establishes that the field candidates are not arbitrary additions to the thesis, but it is neither a claim about every public skill worldwide nor evidence that presence alone improves retrieval.

The controlled experiment keeps non-target information shared among three closely related skills and then shows one operational field. For both BM25 and Qwen, Top-1 improves after the field is shown in every field's own suite. Robustness nevertheless depends on the prompt and retriever: paraphrased outputs are much harder for BM25, while implicit boundaries and paraphrased success or workflow requirements are difficult for both retrievers. Because the seven suites contain different authored cases, these results show that each field can help but do not establish a formal strongest-to-weakest ranking. The field values also vary naturally in length and detail, so the estimate concerns the information as written rather than field meaning after length normalisation. Examples/tests provide the negative control: generic examples do not create a stable distinction, whereas rare near-exact matches can attract a selector without proving a distinct capability.

The public original-document removal experiment tests whether the same information remains consequential when it is embedded in natural artifacts. It compares intact public source documents with the same documents after every cited line for one field, or a defined field group, has been removed from all candidates. Use condition and success/verification have stable strict Hit@1 losses under both BM25 and Qwen. The task-specification and execution/verification groups also have stable losses under both retrievers. Input/precondition, output/artifact, workflow/procedure, dependency/resource, and applicability/capability are positive only under BM25; boundary/not-for has no stable positive result. Thus public documents carry both useful and redundant evidence: some deletions change the strict winner, whereas others leave enough information elsewhere for a selector to recover. The result is deliberately narrower than claiming one unique human reason for every selection, because the edited documents can be incomplete and one operational fact can occur in several sections.

The historical frozen-v0.4 representation matrix motivates, but does not answer, the revised RQ2. It often changes content, structure, candidate generation, and reranking together. For example, with Qwen held fixed, R2 structured representation improves top-1 over R1 flat metadata, and full skill text behaves differently again. These rows show that representation and selector choice interact, but they cannot establish whether the gain comes from additional facts, explicit labels, reduced dilution, or a field-sensitive selector.

The historical public-gold results complicate the story further. On public-gold, R2 is best for Qwen embedding-only retrieval, but R1 plus Qwen reranking has the highest strict top-1. This likely happens because public skills often contain strong provider, platform, or task cues in names and descriptions. In those cases, flat descriptions are not necessarily weak metadata.

\section{RQ2a Answer: Facts Matter More Consistently Than Headings}

RQ2a removes the attribution ambiguities in the historical matrix by holding reviewed propositions and candidate identities fixed. Its strongest conclusion is the operational-content positive control: fielded representations rise by 42.4, 53.1, and 61.8 percentage points over shared-only for Qwen, BM25, and SkillRouter respectively. A selector cannot resolve the near-neighbour decision when candidate-specific operational facts are absent, and all three architectures can exploit those facts when they are restored.

The representation-form result is more qualified. Explicit headings improve BM25 by 2.1 points over matched flat text, which is statistically detectable but below the frozen three-point practical threshold. The same headings reduce Qwen single-vector top-1 by 6.7 points and have no clear effect on SkillRouter cross-encoder selection. This rejects the simple hypothesis that visible field organisation is universally beneficial. Field labels interact with the scoring architecture: they can provide useful lexical anchors, consume or distort a pooled semantic vector, or be largely neutral when a cross-encoder jointly attends to the query and candidate.

The semantic field-aware result sharpens this interpretation. Frozen uniform-top-two scoring improves Qwen fielded selection by 6.6 points, but its 0.823 top-1 is effectively the same as Qwen flat at 0.824. The field-aware mechanism therefore recovers the penalty of compressing labelled fielded text into one vector; it does not beat the best matched Qwen serialisation. Descriptive field slices further show that the global rule helps success/verification, boundary, workflow, and use-condition cases while harming dependency/resource and output/artifact cases. The appropriate conclusion is not ``structure wins'', but that field-wise use is selectively valuable and a single aggregation policy cannot treat all operational fields as interchangeable.

\section{RQ2b Answer: Full-Library Accuracy And Compression Trade Off}

RQ2b tests a different question from RQ2a. It does not hold candidate content fixed, and it does not assign a causal effect to a particular field. Instead, it asks what happens when native metadata, original full bodies, extracted unlabelled evidence, and field-labelled I3C are used to retrieve from the same 2,433-skill library. The strict endpoint contains 381 intended gold labels. It is a demanding full-library retrieval test, but a strict-label miss is not equivalent to downstream task failure or a proof that the returned alternative is useless. V3 is a prospectively frozen internal re-evaluation of historically explored prompts, so its bootstrap intervals characterise fixed-benchmark stability rather than untouched confirmatory generalisation.

The main result is an accuracy--compression trade-off, not a universal I3C benefit. On public-gold B1 retrieval, I2 full bodies outperform I3C for all three first-stage retrievers: 0.644 versus 0.504 for BM25, 0.686 versus 0.538 for Qwen, and 0.806 versus 0.746 for SkillRouter. I3C simultaneously reduces the selector-visible corpus by 72--76\% and recorded document/index time by 64--79\% relative to I2. Fixed-Top-20 reranking changes candidate order but does not consistently reverse this pattern: I3C remains below I2 in five of six public retriever--reranker combinations. The failure decomposition makes the pipeline boundary visible: I3C has no lower candidate-miss rate than I2 in any comparison, and a reranker cannot promote a gold that failed to enter its fixed Top-20 list.

B4 prevents an overly convenient explanation of this result. Only 4 of 755 retained I3 items associated with the strict-gold prompt records are visible in the corresponding I1 descriptions after whitespace folding. I1's competitive result is therefore not explained by literal duplication of the retained I3 operational items in short metadata. Coarse task/category cues, benchmark properties, or changes introduced by evidence selection and serialisation are plausible alternatives, but the present post-hoc census does not causally choose among them. The thesis should not claim that flat metadata is generally sufficient, nor that I3C is generally inferior; it reports the observed V3 trade-off and the limits of the explanation.

\section{Qwen and SkillRouter Represent Different Accuracy--Cost Points}

Although both model families operate on text, the tested conditions are not equivalent embedding systems. Qwen \texttt{text-embedding-v4} is a bi-encoder: it embeds queries and candidate representations independently, permits document-vector precomputation, and ranks by cosine similarity. SkillRouter RQ2a uses \texttt{SkillRouter-Reranker-0.6B}, a cross-encoder that jointly scores each query--candidate pair. It is therefore able to model direct token-level interactions, but the final relevance score cannot be precomputed from the candidate alone.

This distinction explains the observed trade-off. SkillRouter reaches 0.955 top-1 on flat text, the best condition in the matrix, while Qwen flat reaches 0.824. Qwen's document construction is an amortisable one-time cost and its online scoring is about 2.87 ms per prompt after vectors exist. SkillRouter requires about 11.10 s per novel prompt for three candidate pairs in the measured environment. Warm pair-score cache hits are useful for verification but are not a realistic substitute for novel-query inference. RQ2a did not test a SkillRouter-native field-aware architecture; the field-aware row uses Qwen component embeddings, while historical SkillRouter-first-stage plus M6 rows are a different two-stage pipeline.

RQ2b uses the same model names in a different pipeline role. Its B1 SkillRouter condition is a native bi-encoder, comparable to Qwen as a cached document-retrieval stage; its B2 SkillRouter condition is a cross-encoder reranker over the persisted Top-20 list. Qwen B2 is an independent reranker over the same fixed lists. The results should therefore not be collapsed into a statement that one ``SkillRouter score'' or one ``Qwen score'' is universally better: whether a model embeds a full library, ranks a fixed three-sibling set, or reranks Top-20 candidates determines both what it can recover and what cost it incurs.

\section{Full Skill Text And Extracted I3C Have Different Failure Modes}

The RQ2b result rules out two simplistic prescriptions: ``always embed the whole document'' and ``always extract I3C''. Original full bodies retain detail that helps candidate inclusion on the public-gold endpoint, while I3C is a much smaller selector surface. The best choice in this study depends on whether strict top-1 and Recall@20 or document/index burden are the more binding deployment constraint.

The historical frozen-v0.4 and external SkillRouter-Eval-Core rows remain useful as portability context, but they use different prompts, source artifacts, metrics, and extraction versions. They cannot override the source-grounded V3 full-library comparison. In particular, an external BM25 result where an earlier I3C extraction is favourable does not establish that I3C must be favourable for every library or first-stage retriever.

\section{Lexical Baselines Are Stronger Than Expected}

The historical M6-v1 local runs show that lexical candidate generation remains a serious baseline. Controlled TF-IDF/schema plus M6-v1 field-aware reranking reaches 69.8\% top-1 and 97.1\% top-5. Public-gold TF-IDF/flat plus task-field reranking reaches 68.1\% strict top-1.

This result changes the framing of the thesis. The project should not be presented as a simple contest between dense embeddings and structured representations. A more accurate framing is:

\begin{quote}
Which operational selection information should be exposed, and which retrieval or reranking strategy can use it reliably under semantic confusability?
\end{quote}

In the current benchmark, lexical retrieval can exploit strong names, descriptions, output words, and workflow terms. That does not make the benchmark invalid, because real systems also contain lexical cues. But it means the final analysis must treat lexical+field-aware methods as serious baselines and investigate whether their wins come from genuine procedural matching or from surface cue overlap.

\section{What Information Appears Useful}

The controlled field-isolation suites suggest that output artifacts, input conditions, preconditions, workflow steps, procedural cues, use conditions, boundary evidence, dependency/resource constraints, and concrete success criteria can all function as retrieval evidence. These fields encode procedural suitability: what the skill is for, what it expects, what it produces, how it operates, what scope or authority limits apply, what external capability it requires, and how completion is judged.

The use-condition result is the most direct version of this claim because the field states when the skill should be selected. It is clearly discriminative in the controlled clusters, but the paraphrase rows show that it is not merely a perfect direct hint. BM25 drops more sharply under paraphrase than Qwen embedding, which suggests that task/use-condition matching sometimes requires semantic interpretation of the user's intent rather than only lexical overlap with the skill's trigger wording.

The boundary/not-for result also shows why not all fields should be treated as the same kind of evidence. Explicit boundary prompts produce strong gains, but implicit-authority prompts are much harder. When the user asks for an internal review note, triage packet, or evidence summary, a human may infer the authority limit, but ordinary lexical and embedding retrievers often do not. Boundary information is therefore better understood as a conditional scope or guardrail signal than as a simple positive matching field.

The dependency/resource result gives a second example of conditional evidence. In controlled clusters, capability compatibility is highly discriminative: when the skill-side field states a required platform, tool, API, runtime, permission, version, or resource, both BM25 and Qwen embedding can use it to break near-neighbour ties. The difficult part is acquisition and normalisation. A user may not say ``React 18'', ``Chrome DevTools MCP'', or ``GitHub token'' in the natural request. The router may need to obtain that cue from repository files, configured tools, credentials, package metadata, or runtime state. Public skills also mix hard requirements with generic setup text: \texttt{npm install}, \texttt{npm ci}, or ordinary build/test commands are usually execution preparation rather than identity evidence, while Chrome DevTools MCP, \texttt{GH\_TOKEN}, or Node.js 20+ are capability constraints. Therefore the thesis should not claim that raw dependency prose is always useful. It should claim that specific external capability compatibility is useful when the selector has access to sufficient capability context, whether that context comes from the request or from the routing environment.

Success/verification is a third conditional field type. It is useful when the request names the acceptance gate directly, but it is much less robust under paraphrase: in the completed suite, both BM25 and Qwen embedding fall to 48.0\% top-1 on paraphrased success criteria. This does not make the field useless. It shows that success criteria are often semantically close to one another, such as schema validation, evidence anchoring, threshold checks, and smoke checks. The field is therefore better interpreted as a completion-gate or late-stage tie-breaker signal than as a strong first-pass discriminator.

The public-gold and public-audit evidence is more cautious. Public descriptions can already contain strong provider, platform, task, or artifact cues, and public skill bodies often mix useful routing evidence with setup text, examples, broad parent instructions, and duplicated capabilities. This does not mean operational fields are useless. It means the controlled suites provide the clearest evidence that exposing a field can help, while the public experiment tests whether removing the same kinds of information matters in less controlled documents.

The most defensible current claim is therefore not that every field always improves retrieval. The current claim should be:

\begin{quote}
Candidate-specific operational information can improve routing when it is visible and matches the request, but the size and reliability of that improvement depend on the field, prompt, retriever, and surrounding document.
\end{quote}

RQ1 is not a competing retrieval architecture. SSL proposes a complete layered skill representation, while SkillRouter and SkillRet evaluate larger retrieval pipelines \cite{liang2026ssl,zheng2026skillrouter,cho2026skillret}. RQ1 asks a narrower question: when several skills are plausible, which kinds of operational information can change the selection? The controlled addition and public removal experiments test that question separately from architecture. They therefore complement these systems rather than claiming to outperform them; representation and retrieval architecture are compared under RQ2.

\section{Why Existing Field-Aware Prototypes Are Diagnostic}

M6-v1-local is useful because it operationalises the thesis claim: it parses request-side cues and compares them with skill-side fields. It also separates first-stage candidate recall from reranking quality. This helps diagnose whether failures come from the retriever excluding the gold skill or from the reranker ordering candidates poorly.

But M6-v1-local is not yet a final robust structure-aware method. It is still largely lexical. It can miss semantically equivalent wording, overvalue repeated field terms, and under-handle contradiction, negation, hierarchy, and broad parent skills. Its public-gold weakness is therefore informative. It shows that field-aware retrieval needs semantic matching or LLM-assisted field interpretation, not only manually designed field names.

M6-v2 already provides a semantic field-matching prototype, but it remains tied to the old frozen-v0.4 representations, first-stage blending, selected field sets, request parsing, penalties, and calibration. The completed RQ2a adapter reuses M6-v2's embedding and caching components while removing first-stage blending, adding all seven fields, hiding target-field metadata, and freezing one target-agnostic aggregation rule on development clusters. Graph, tree, and DAG variants are outside the completed RQ2a scope.

\section{Benchmark Validity}

The benchmark has several strengths. It contains controlled confusable clusters, a large 2433-skill scale condition, generated background skills, imported public skills, public-gold prompts, acceptable alternatives, and validation checks for integrity, procedural distinctness, semantic confusability, and leakage.

The benchmark also has limitations. The controlled skills were created with AI assistance under a researcher-defined rubric and then checked and approved by the researcher. They may therefore express procedural distinctions more cleanly than publicly available skills, and some gold field values are longer or more detailed than their alternatives. The public experiment uses purposefully selected, removal-eligible groups rather than a random sample. It tests less controlled documents, but its result is still limited by group selection, field-specific eligibility, broad parent skills, duplicate capabilities, provider-specific names, and confirmation of the frozen gold labels by the same researcher with the labels visible.

The external SkillRouter-Eval-Core benchmark has a different limitation. It is much larger and externally sourced, but its scored tasks are more informative than the local controlled prompts. The median scored task is 169 rough tokens, and most tasks include file/path references plus explicit input and output cues. Its skill bodies are also highly markdown-normalised, with frequent headings such as \texttt{overview}, \texttt{when to use}, \texttt{workflow}, \texttt{quick start}, and \texttt{dependencies}. This can make full-body retrieval more attractive than it would be in a shorter, near-neighbour controlled setting. The external benchmark therefore strengthens portability evidence, but it does not remove the need for the local controlled benchmark.

The benchmark should therefore be described as a stress test for semantic confusability, not as a universal public benchmark. Its value is that it isolates a specific failure mode: skills that look semantically plausible but differ procedurally. Its limitation is that final generalisation requires more public-skill validation and downstream testing.

\section{Failure Modes}

The current experiments reveal several recurring failure modes:

\begin{itemize}
  \item \textbf{First-stage exclusion:} the gold skill is not included in the candidate set, so reranking cannot recover it.
  \item \textbf{Object-procedure confusion:} the selector focuses on the shared task object, such as PDF, GitHub, or spreadsheet, rather than the required procedure.
  \item \textbf{Full-document dilution:} full skill text contains the right information but also enough unrelated text to weaken embedding similarity.
  \item \textbf{Provider cue dominance:} public skill names or providers become strong cues, sometimes correctly and sometimes superficially.
  \item \textbf{Boundary failure:} negative or avoid conditions are treated as ordinary positive text or ignored.
  \item \textbf{Broad parent selection:} a broad public skill is selected instead of a more specific atomic skill.
  \item \textbf{Lexical overfitting:} a method ranks the correct skill because of surface term overlap rather than deeper procedural compatibility.
  \item \textbf{Underspecification:} the user request lacks enough evidence to choose a stable skill without clarification.
\end{itemize}

\section{Implications for Skill Retrieval Systems}

The practical implication is that scalable skill libraries should not rely only on names and short descriptions, but they also should not assume that full-document embedding solves the problem. A better architecture should separate first-stage candidate generation from final selection, preserve procedural selection information, and use field-aware logic for outputs, workflows, preconditions, dependencies, and boundaries.

The current results also suggest that multiple retrieval strategies may be economically useful. Lexical retrieval is cheap and strong when skill names and descriptions are well written. Dense retrieval is useful for broader semantic recall but can be sensitive to representation length and noise. Reranking improves ordering, but its value depends on whether the candidate set contains the gold skill and whether the reranker sees the right information.

Cost should therefore be analysed as an amortised pipeline rather than as a single per-query number. RQ2b measures the selector-visible corpus, B1 document/index work, B1 query work, and fixed-Top-20 reranking work. It does not contain a durable comparable I3C extraction wall-time or provider invoice, so it cannot calculate a true break-even query count. The appropriate claim is narrower: I3C lowers recorded B1 representation/index burden relative to I2, while the missing construction-cost term and the observed strict-label loss must both remain visible in any deployment decision.

\section{Current Limitations}

RQ2a is complete, but its scope is deliberately narrow. It is a fixed-candidate mechanism study in which the gold skill is always present among three controlled siblings. RQ2b complements it with full-library candidate generation, Recall@20, fixed-Top-20 reranking, and a public-gold stratum, but it cannot isolate the effect of one field because the representations differ in content and extraction. In the controlled suites, the candidate text is identical before the target field is added, so the observed change measures the effect of showing that field content in that suite. Because the field values were not matched for length or detail, the experiment does not separate field meaning from wording or amount of information. The V3 public and controlled full-library sources provide strict-label retrieval evidence under a larger, messier candidate pool.

The selector coverage also bounds the claim. The Qwen field-aware condition uses one frozen uniform-top-two rule; field-level slices are descriptive and cannot justify post-hoc field-specific weights. No SkillRouter-native field-aware condition was tested. The RQ2b availability census is post-hoc and provenance-bounded: it shows where retained evidence occurs, not which information caused a retrieval result. Its strict labels are intended routing targets rather than an acceptable-alternative or downstream-success endpoint. Automatic source/provenance validation does not establish independently reviewed semantic completeness of I3C: a protocol-invalid manual review was quarantined and no replacement semantic-completeness audit was used. The cost measurements are implementation- and hardware-specific, and RQ2b lacks a complete I3C construction ledger. M4 tree routing, M5 graph retrieval, full-scale progressive disclosure, and downstream task success remain future or separate studies rather than missing rows from the completed RQ2a/RQ2b evidence.
